% =====================================================================
%  Accelerating the Flux MMDiT on PLENA  (benchmarked on Open-Sora)
%  BEng final report — Imperial College London EEE
%
%  STRUCTURE-ONLY SKELETON.
%  This file currently lists ONLY the chapter / section structure and a
%  one-line note on what each part should contain. No body text yet —
%  content is filled section-by-section on request.
%
%  Author: Ziqian Gao
% =====================================================================

% ---------------------------------------------------------------------
% FRONT MATTER (handled by the Imperial template main.tex / preamble):
%   - Title page  (\title, \author=Ziqian Gao, \cid, \supervisor, \degree=BEng, \course, \submityear=2026)
%   - Originality statement & copyright
%   - Acknowledgements
%   - Abstract                  → chapters/Abstract.tex
%   - Table of contents / list of figures / list of acronyms
% ---------------------------------------------------------------------


% =====================================================================
% ABSTRACT
%   what:  one paragraph. We accelerate the Flux MMDiT architecture
%          (double- + single-stream blocks) on the PLENA accelerator,
%          using Open-Sora as the benchmark model. Analyse three things:
%          accuracy (quantisation), latency, power. Headline numbers.
% =====================================================================


% =====================================================================
% CHAPTER 1 — INTRODUCTION
% =====================================================================
%   1.1 Motivation
%       what: video/image diffusion is expensive; the MMDiT backbone
%             dominates cost; dedicated accelerators (PLENA) vs A100.
%   1.2 The Flux MMDiT and Open-Sora as benchmark
%       what: Flux introduced the MMDiT (double-stream + single-stream);
%             Open-Sora-v2 reuses it — we use its real dims/weights as
%             the benchmark to measure on.
%   1.3 Objectives
%       what: the three analysis axes — accuracy, latency, power —
%             across inference (low precision) and training (high prec).
%   1.4 Contributions
%       what: bullet list of what this report delivers.
%   1.5 Report structure
%       what: one paragraph road-map of the chapters.


% =====================================================================
% CHAPTER 2 — BACKGROUND
% =====================================================================
%   2.1 Diffusion transformers and the MMDiT
%       what: DiT / MMDiT basics; double-stream vs single-stream block;
%             joint attention; RoPE, QK-norm, gelu, modulation, residual.
%   2.2 The Flux / Open-Sora block structures
%       what: SSB and DSB topology in detail (the two block types we target).
%   2.3 Quantisation for inference (MX-E4M3 / FP8)
%       what: low-precision number formats, per-channel/per-token scales,
%             accumulation precision (FP8 -> FP16/FP32 accumulate).
%   2.4 The PLENA accelerator
%       what: systolic matrix core (BLEN x MLEN), vector unit, SRAMs,
%             HBM, the ISA at a high level; device-level parallelism.
%   2.5 The TileLang compilation flow   <<< COMPILER DEMOTED TO HERE
%       what: ONE short section. We write kernels in ordinary GPU
%             TileLang style and a compiler lowers them to PLENA ISA.
%             NO pass-by-pass detail — just "this is the tool we use".


% =====================================================================
% CHAPTER 3 — ANALYSIS AND DESIGN
% =====================================================================
%   3.1 Where the compute is
%       what: FLOP breakdown — MMDiT dominates the model; SSB/DSB
%             dominate the MMDiT pass; pick the targets.
%   3.2 Precision design: inference (low) vs training (high)
%       what: MX-E4M3 forward vs fp32 backward; the power model's two
%             legs (MCU by accumulation precision, vector/SRAM by bit-width).
%   3.3 Matching the hardware to an A100 (the comparison budget)
%       what: count MACs both sides; pick 12 BLEN-8 devices ≈ one A100.
%   3.4 Device allocation
%       what: split one kernel across N devices at the mid-IR boundary
%             (blockIdx per device + block-id loop for the remainder).


% =====================================================================
% CHAPTER 4 — IMPLEMENTATION
% =====================================================================
%   4.1 Forward block implementation
%       what: SSB/DSB forward in TileLang at real dims; kernel set.
%   4.2 Backward (training) implementation
%       what: per-operator backward; transpose-A buffer (MRAM-matched);
%             linear dX/dW as plain matmuls; flash-attention backward.
%   4.3 Optimiser step (AdamW)
%       what: vector-only AdamW; element-wise update.
%   4.4 The verification methodology (sim → GPU → ISA)
%       what: simulator anchors numerics on a layer; reproduce the same
%             quantisation on GPU for full-model accuracy; ISA analysis
%             gives latency/energy. (the 3-stage chain)


% =====================================================================
% CHAPTER 5 — RESULTS AND EVALUATION
% =====================================================================
%   5.1 Accuracy
%       what: forward numerical correctness (TileLang vs PyTorch);
%             FP8 quantisation (W8A16/W8A8, PTQ vs QAT); error propagation
%             over depth + denoising loop; whole-MMDiT final precision.
%   5.2 Latency and energy — inference
%       what: per-block forward vs A100; full MMDiT pass; full generation.
%   5.3 Latency and energy — training
%       what: backward per-block (low/high precision) vs A100; AdamW tail.
%   5.4 Power-model validation
%       what: bottom-up vs DC-synthesis cross-check (the George comparison).


% =====================================================================
% CHAPTER 6 — REFLECTIONS  (template-required, IET sections)
% =====================================================================
%   6.1 Legal and Ethical matters
%   6.2 Environmental and Social impact
%   6.3 Equality, Diversity, and Inclusion
%   6.4 Quality Management Systems


% =====================================================================
% CONCLUSIONS AND FUTURE DIRECTIONS
%   what: three-axis summary (accuracy near-lossless FP8; latency ~A100;
%         energy advantage); limitations; future work.
% =====================================================================


% =====================================================================
% APPENDICES
%   A.  End-to-end A100 validation of the analytic generation numbers
%       (the 0.5%-error table).
%   B.  Power-model details (unit-power derivation, precision scaling).
% =====================================================================
